\documentclass{article}

\usepackage{graphicx}
\usepackage[utf8]{inputenc}
\usepackage[italian]{babel}
\usepackage{fancyhdr}
\usepackage{hyperref} % Pacchetto per gli url e link nel documento
\usepackage{longtable}
\usepackage[table]{xcolor}
\usepackage{makecell}
\usepackage{geometry}
\usepackage{array}
\usepackage{lastpage}
\usepackage{float}


% Definisce un nuovo tipo di colonna per le tabelle. E' come p{}, ma con il testo all'interno della cella già centrato.
\newcolumntype{P}[1]{>{\centering\arraybackslash}p{#1}}

% Setup pacchetto "hyperref"
\hypersetup{
	colorlinks=true,
	linkcolor=black,
	urlcolor=blue
}

\setlength{\headheight}{26.87454pt}
\addtolength{\topmargin}{-14.87454pt}

\usepackage{enumitem}

% Definisce i margini e il formato delle pagine
\geometry{a4paper, top=3cm, left=2.5cm, bottom=3cm, right=2.5cm}

% AGGIUNTA PARAGRAFI E SOTTOPARAGRAFI
\setcounter{secnumdepth}{5}
\setcounter{tocdepth}{5}

% INDENTAZIONE SOTTOPARAGRAFI
\makeatletter
\renewcommand\subparagraph{
\@startsection{subparagraph}{5}{0mm}{-\baselineskip}{.5\baselineskip}{\normalfont \normalsize \bfseries }}
\makeatother

% INDENTAZIONE PARAGRAFI
\makeatletter
\renewcommand\paragraph{
\@startsection{paragraph}{4}{0mm}{-\baselineskip}{.5\baselineskip}{\normalfont \normalsize \bfseries }}
\makeatother

\title{SecondRead}
\date{2021/2022}

\begin{document}
\pagenumbering{gobble}

\makeatletter
    \begin{titlepage}
        \begin{center}
            \includegraphics[width=0.4\linewidth]{images/icon.jpg}\\[10ex]
            {\huge \bfseries  \@title }\\[5ex]
            {\large \@date} \\[10ex]
            {\large \textit{Componente:} Nicholas Pilotto - 1230237} \\[2ex]
            {\large \textit{Componente:} Michele Bonavigo - 1216752} \\[2ex]
            {\large \textit{Componente:} Mattia Casarotto - 1221948} \\[2ex]
            {\large \textit{Componente:} Anita Meta - 1192890} \\[2ex]
            {\large \textit{Destinatari:} Prof.ssa\ Ombretta Gaggi} \\[2ex]
        \end{center}
    \end{titlepage}
\makeatother
\thispagestyle{empty}
\newpage

% Head pagine dopo la prima (con stile fancy)
\pagestyle{fancy}
\fancyhf{}
\rhead{\textbf{SecondRead}}
\lhead{\includegraphics[width=1.1cm]{images/icon.jpg}}

\setcounter{table}{0} % Resetta il contatore delle tabelle

\tableofcontents \pagebreak

% Foot pagine dopo gli indici
\pagenumbering{arabic}
\rfoot{Pagina \thepage/\pageref*{LastPage}}

\newpage
% ============================= Abstract ==============================
\begin{abstract}
Det här avsnittet ska helt enkelt vara just detta: en sammanfattning av hela rapporten. En lämplig omfattning är c:a 200 – 250 ord.  En bra tumregel är att sammanfattningen ska hållas så kort det går, den ska vara kompakt men fortfarande tydlig, informativ och väcka intresse. Ge de viktigaste fakta
och summera allt det som är väsentligt i rapporten.  Följande bör ingå:
\begin{itemize}
\item[--]	Presentation/introduktion av området för arbetet
\item[--]	översiktlig presentation av uppgiften inklusive syfte och frågeställning
\item[--]	Motivation till varför området och uppgiften är viktiga och intressanta
\item[--]	Generell beskrivning av hur du angripit uppgiften, vad du har gjort
\item[--]	Sammanfattning av resultat och slutsatser och vad ditt arbete bidrar med
\end{itemize}

Inga detaljer ska vara med i sammanfattningen, inte heller beskrivning av hur rapporten är uppställd. 
Sammanfattningen ska kunna läsas helt fristående från resten av rapporten, och av en ganska bred grupp av läsare. Den ska ge en bra grund för att en läsare ska kunna bedöma om hen är intresserad av att läsa hela rapporten. 
Sammanfattningen är den del av en rapport som läses allra mest och av flest personer. Därför är det extra viktigt att du skriver en bra sammanfattning. Du behöver ha ett ordentligt grepp om innehållet i rapporten när du skriver sammanfattningen, och när hela rapporten är klar bör du granska och vid behov revidera sammanfattningen så att den överensstämmer med rapporten.

\end{abstract}
\newpage
\section{Analisi}

\subsection{Analisi delle caratteristiche utenti}
\textit{SecondRead} rappresenta una piattaforma intermediaria per la vendita dei libri tra individui e le grandi case editrici. L’utenza primaria del sito sarà quindi composta da lettori, siano essi occasionali o appassionati.

Questo non esclude tuttavia il fatto che anche un soggetto legato all’ambito accademico come un professore oppure uno studente possa utilizzare la pagina per trovare libri riguardanti il proprio corso di studio.

Cercando di includere ogni ambito di interesse per i lettori, sono state create inizialmente circa 12 diverse categorie di libri. Si cerca sempre di avere una quantità sufficiente di copie disponibili per ogni libro in modo da poter offrire agli utenti del sito un servizio impeccabile.

Il sito offre all’utente la possibilità di svolgere una ricerca veloce e precisa, effettuata di solito dalle persone che sanno già cosa vogliono, ma anche la possibilità di svolgere ricerche ad ampio spettro che riguardano argomenti correlati a quello inizialmente cercato.

Un altro obiettivo è stato quello di includere nel \textit{target group} tutte le diverse fasce d’età, motivo per il quale si è optato per un linguaggio semplice, informale e intuitivo che possa essere comprensibile a tutti.

La stessa filosofia è stata applicata anche per quanto riguarda struttura e layout del sito: esso è veloce, intuitivo e responsive in modo da adattarsi a vari tipi di dispositivi. Tutto questo rispettando le principali convenzioni del web.

\subsection{Possibile ricerche sui motori di ricerca}
Elencate di seguito, in ordine di rilevanza (da particolare al generale), le possibili ricerche che dovrebbero presentare tra i risultati il sito web.

La ricerca a cui deve sicuramente rispondere il sito è quella contenente il nome stesso, ossia \textit{SecondRead}. Queste ricerche saranno maggiormente effettuate da utenti del sito, da persone che lo hanno visitato precedentemente oppure da persone a cui è stato raccomandato o riferito, ad esempio da un amico.

Essendo un sito di libri online, è fondamentale che risponda in una ricerca a tutte le \textit{query} che contengono un \textit{ISBN, un titolo di libro, oppure un nome di autore o casa editrice} dei libri che sono presenti nella pagina. Questo tipo di ricerca viene effettuata dalle persone che stanno cercando un particolare libro come per esempio uno studente, un genitore o anche un professore universitario.

È importante che il sito risponda anche alle ricerche ad ampio spettro. Ad esempio, se un appassionato di scienza cerca “Scienza e Fantascienza” oppure un appassionato dei manga cerca \textit{“Fumetti e Manga”}, essendo queste categorie di libri parte del sito, esso dovrebbe apparire tra i risultati della ricerca.

Uno degli obiettivi principali del team è quello di espandere la clientela ed attirare nuovi utenti. In modo da realizzarlo, è necessario che il sito risponda anche alle ricerche generali che possono contenere parole chiave come: \textit{libri, libro online, vendita libri, dove comprare libri online,..} . Queste sono le ricerche che vengono effettuate di solito da persone che non hanno esperienza e preferenze specifiche ma che vogliono iniziare a leggere oppure hanno recentemente sviluppato la lettura come \textit{hobby}.

Un altro caso che aiuterebbe ad aumentare la visibilità, è quello in cui la ricerca contiene parole chiave come: libri bestseller, libri scontati online, nuove uscite, prezzo sotto 5 euro...

Una persona che vorrebbe fare un regalo senza spendere troppo sarebbe il caso di utilizzo di questo tipo di ricerca.

In generale, essendo il sito progettato in modo da poter essere utilizzato da tutte le fasce d'età, si è prestato attenzione alle ricerche effettuabili da utenti appassionati ed esperti, ma anche a quelle effettuate da principianti.

\subsection{Conclusioni}
Partendo dall’analisi e dalle considerazioni fatte precedentemente riguardanti il sito e gli utenti, otteniamo le seguenti conclusioni:
\begin{itemize}
	\item \textit{Utenza finale}: gli appassionati di libri e persone legate all'ambito 	accademico
	\item \textit{Target group}: tutte le varie fasce d’età
	\item Ricerca: è possibile effettuare vari tipi di ricerca. Nel caso di ricerca
			specifica, si può utilizzare direttamente la barra di ricerca, mentre per
			quella ad ampio spettro si possono trovare informazioni nei raggruppamenti di
			generi o fasce di prezzo
	\item \textit{Interazione}: all’interno del sito è possibile interagire in diversi
			modi. Un esempio di
			inserimento di dati si trova nella creazione di un account oppure 
			nell'aggiunta di un indirizzo. Un esempio di cancellazione di dati lo 
			proviamo nella cancellazione di recensioni memorizzate precedentemente.
			Si ha anche la possibilità di modificare i dati inseriti come la password,
			l’username oppure l’email.
	\item \textit{Community}: le recensioni sono un elemento pensato ed aggiunto al sito 	in modo da poter creare una forma solida di comunicazione ed interazione tra 	gli utenti. Leggendo le diverse opinioni presenti sulla pagina di un
			determinato libro, l’utente crea un'idea generale più chiara del libro e
			decide se acquistarlo.
\end{itemize}

\subsection{Utenti}
Abbiamo 3 diverse tipologie di utenti:
\begin{enumerate}
	\item \textit{Utente generico(generale)}: Un utente viene considerato generico quando:
		\begin{itemize}
			\item Non è stato registrato, ossia non ha creato il suo account sul sito.
			\item Possiede un account ma non ha ancora effettuato l’accesso.
		\end{itemize}
		Questa tipologia di utente ha accesso a buona parte del sito ma subisce restrizioni su alcuni servizi offerti. Un utente generico non può, ad esempio, fare ordini oppure lasciare recensioni. Tuttavia è comunque in grado di:
			\begin{enumerate}[label*=\arabic*.]
				\item Visualizzare la Homepage del sito
				\item Visualizzare la pagina “Ricerca”
				\item Visualizzare la pagina “Generi”
				\item Visualizzare la pagina “Bestseller”
				\item Visualizzare la pagina “Offerte”
				\item Visualizzare l’elenco dei libri listati in ogni categoria
				\item Cercare un libro nella pagina di ricerca
				\item Visualizzare tutti i dettagli di un particolare libro in vendita 
				\item Registrarsi
				\item Effettuare l’accesso
			\end{enumerate}
	\item \textit{Utente loggato}: Rappresenta l’utente che ha effettuato l’accesso al suo account personale. Eredita tutti i casi d’uso dell’utente generico tranne la parte di registrazione e dell’accesso. Dispone di altre funzionalità aggiuntive come:
			\begin{enumerate}[label*=\arabic*.]
				\item Visitare l'Area Privata 
				\item Visualizzare i suoi ordini
				\item Gestire gli indirizzi
				\item Gestire le recensioni
				\item Visualizzare e modificare le informazioni di login 
				\item Gestire ed aggiornare la propria \textit{wishlist}
				\item Effettuare un ordine
				\item Lasciare una recensione
				\item Contattare il numero di assistenza in caso di problemi
			\end{enumerate}
	\item \textit{Amministratore(admin)}: Rappresenta l’utente che gestisce il sito ed è identificato dall’indirizzo email: \underline{admin@gmail.com}. Possiede la maggior parte delle funzionalità degli utenti sopra elencati ma con alcune eccezioni. Non può ad esempio:
		\begin{itemize}
			\item Lasciare una recensione
			\item Aggiungere elementi nella \textit{wishlist} oppure nel carrello
			\item Effettuare un ordine
		\end{itemize}
Dopo aver effettuato l’accesso al suo account, attraverso l’Area Privata sulla navbar, può eseguire le seguenti funzionalità:
		\begin{enumerate}[label*=\arabic*.]
			\item Aggiungere nuovi libri
			\item Modificare dettagli di un libro già presente nella pagina 
			\item Visualizzare i dati analitici del sito
		\end{enumerate}
\end{enumerate}
\newpage
\section{Struttura del sito}\label{sec:struttura}

\subsection{Header}
Il \textit{header} è sempre lo stesso ed è presente su tutte le pagine. Contiene il logo e il nome del sito in alto a sinistra.

\subsection{Breadcrumb}
La \textit{breadcrumb} è uno strumento fondamentale per quanto riguarda l’accessibilità. La funzione principale è quella di aiutare l'utente ad orientarsi sulla propria posizione all'interno della pagina \textit{web}. Si trova immediatamente sotto il menu ed è presente su tutte le pagine, sia nella versione \textit{desktop} sia in quella \textit{mobile}.

Contiene i \textit{link} relativi alle pagine che sono state visitate precedentemente dall’utente. L'ultima pagina listata è quella corrente che non contiene un \textit{link} in modo da evitare una circolazione non necessaria.

\subsection{NavBar}
La \textit{navbar} costituisce una delle parti più importanti del sito siccome contiene i vari link per la navigazione. Nella maggior parte delle pagine, \textit{header} e \textit{navbar} sono invariati ma la cosa cambia in base al tipo di utente.
	\begin{enumerate}
		\item Utente generico: nel caso in cui l’utente non possiede un \textit{account} oppure non ha ancora effettuato l’accesso, la \textit{navbar} è costituita dalle pagine:
			\begin{itemize}
				\item Ricerca
				\item Generi
				\item Bestseller
				\item Offerte
				\item Area Riservata
			\end{itemize}
		Area Riservata non è in sé una pagina come le altre, ma serve a creare il collegamento con la pagina Accedi. L’utente generico, dopo aver cliccato su Area Riservata, viene indirizzato alla pagina Accedi dalla quale può direttamente effettuare l'accesso all'\textit{account}.

		In alternativa, se l’utente generico non possiede un’account \textit{SecondRead} (non lo ha creato precedentemente), può farlo direttamente cliccando il pulsante Registrati presente nella pagina Accedi. Verrà indirizzato alla pagina Registrati dove deve compilare il \textit{form} per la registrazione.

		Si è scelto di aggiungere un ulteriore livello di profondità (per andare alla pagina Registrati si deve effettuare un “\textit{click}” in più rispetto a quello che si doveva fare se la pagina fosse stata direttamente presente sulla \textit{navbar}) con lo scopo di evitare il sovraccaricamento del \textit{menu} e soprattutto in modo da evitare il sovraccarico cognitivo dell’utente.

		Allo stesso tempo, attraverso questa organizzazione l'adattamento della \textit{navbar} per i dispositivi con schermi più piccoli risulta più facile.

		\item \textit{Utente loggato}: l’utente ha già effettuato l’accesso sul proprio account. La sua \textit{navbar} è composta dalle pagine:
			\begin{itemize}
				\item Ricerca
				\item Generi
				\item Bestseller
				\item Offerte
				\item Area Riservata
				\item Carrello
				\item Esci
			\end{itemize}
		In questo caso, l’Area Riservata esegue una funzione diversa. Quando l’utente lo clicca, viene indirizzato alla pagina \textit{Account} dove si trovano tutte le sezioni e funzionalità che possiede.
		\item \textit{Amministratore}: in questo caso la \textit{navbar} è leggermente diversa da quella dell'utente loggato siccome non contiene la pagina Carrello. La funzione dell'Area Riservata rimane la stessa, ossia dopo il \textit{login} se cliccata indirizza alla pagina \textit{Account} dell'amministratore (admin).
	\end{enumerate}
	
\subsection{Contenuto della pagina}
Le informazioni più importanti sono direttamente presenti nella home page siccome è la parte del sito che l’utente vede appena arrivato. Abbiamo le sezioni Bestseller, Nuove uscite e A meno di 5 euro.
Si è utilizzato lo schema a tre pannelli che aiuta a dare risposta alle 3 domande importanti:
\begin{enumerate}
	\item \textit{Di cosa si tratta?}: La risposta a questa domanda si trova nel contenuto della pagina home
	\item \textit{Dove sono?}: La prima risposta a questa domanda si trova nel titolo della pagina ma si ottiene
		una risposta anche guardando nel \textit{breadcrumb}
	\item \textit{Dove posso andare?}: La risposta si trova sul menu dove sono elencate tutte le altre pagine che
		è possibile visitare.
\end{enumerate}	
	
	
	
	
	
	
	
	
	
	
	
	
	
	
	
\newpage
\section{Progettazione del sito}\label{sec:progettazionesito}

\subsection{Obiettivi}
La fase di progettazione e sviluppo del sito \textit{SecondRead} è stata fatta tenendo sempre in conto alcuni obiettivi e principi essenziali di un sito ben organizzato e accessibile.
\begin{itemize}
	\item \textit{Separazione contenuto, presentazione e struttura}: Uno degli obiettivi più importanti nella progettazione. 
	
		Il contenuto è stato realizzato tramite documenti \textit{HTML5} e \textit{PHP}. La parte della presentazione grafica è realizzata attraverso i file esterni \textit{CSS} che vengono poi richiamati nelle pagine \textit{HTML}. La realizzazione del comportamento avviene tramite gli \textit{script} \textit{Javascript}. Tutto il codice delle pagine è stato scritto secondo le raccomandazioni \textit{W3C} e in ogni fase della progettazione si è verificato che tutto andasse bene tramite la validazione.
	\item \textit{Accessibilità}: L’accessibilità è uno dei punti fondamentali da tenere in mente durante la progettazione. Si deve realizzare un sito che possa essere utilizzabile da tutte le categorie di utenti e da tutti i tipi di dispositivi e motori di ricerca. In modo da raggiungere questo obiettivo, abbiamo utilizzato varie tecniche.
		\begin{enumerate}
			\item Sono state fornite equivalenti testuali leggibili dallo \textit{screen reader} per tutti i tipi di media come supporto per gli utenti con disabilità visiva.
			\item Grazie all’utilizzo delle misure relative o in percentuale nei fogli \textit{CSS} otteniamo un sito con un design fluido dove ogni pagina è responsive, ossia adattabile ai schermi di diverse dimensioni utilizzati dagli utenti.
			\item Sono stati usati gli attributi \textit{lang} per le parole in lingue diverse dall’italiano in modo da essere lette correttamente dallo \textit{screen reader}.
			\item La scelta dei colori del sito è stata fatta in modo da garantire un contrasto adeguato tra i diversi elementi strutturali e facilitare la lettura del contenuto a chi soffre di disturbi visivi.
			\item Con lo scopo di rendere il sito facile da usare, e cercando di aiutare l’utente alla creazione di una mappa mentale, si sono rispettate alcune convenzioni importanti del web come il colore dei \textit{link}. I \textit{link} non visitati sono sottolineati e in color blu mentre quelli visitati diventano viola.
			\item Si è aggiunto il testo \textit{“Ti trovi in :”} che precede in \textit{breadcrumb} per aiutare gli utenti svantaggiati.
			\item Abbiamo aggiunto un pulsante in basso chiamato \textit{“Torna su”} che se cliccato, effettua uno \textit{scroll} verso l’alto fino all’inizio della pagina, migliorando così l’esperienza durante lo \textit{scroll} delle pagine lunghe.
			\item Utilizzo del \textit{tag abbr} per definire le abbreviazioni.
			\item Si è utilizzato un linguaggio semplice, informale e intuitivo da poter essere
				comprensibile a tutti e in modo da aiutare l’utente a navigare all’interno della pagina corrente oppure tra le diverse pagine.
		\end{enumerate}
	\item \textit{Layout}: Il sito è stato progettato con un layout a tre pannelli:
		\textit{Header}, \textit{Navbar}, Contenuto.
	\item \textit{Resize Mobile}
	\item \textit{Schema organizzativo}
\end{itemize}
\newpage
\section{Implementazione}\label{sec:implementazione}

\subsection{Linguaggi e strumenti}
\subsubsection{HTML}
Il linguaggio usato per la parte della struttura e contenuto del sito è \textit{HTML5}. Questa scelta è stata fatta per due principali motivi:
\begin{enumerate}
	\item Essendo un sito di vendita online, è importante che tutte le pagine siano le più coerenti possibili in modo da poter competere nel mercato con un sito fatto bene ed aggiornato. Dall’altra parte, avere pagine scritte con linguaggi recenti aiuta moltissimo nella leggibilità del codice ed è più facile poter fare modifiche.
	\item Un altro motivo per il quale si è scelto \textit{HTML5} è correlata con l’utenza finale del sito. Sono maggiormente le generazioni più giovani quelle che si affidano al mondo online per fare acquisti. Di conseguenza, saranno i dispositivi e i motori di ricerca più recenti quelli che si useranno di più.
\end{enumerate}
Si è comunque mantenuta la compatibilità con \textit{XHTML}, permettendo al sito di essere pienamente funzionante anche sui browser più vecchi e obsoleti.
Durante tutto il processo di progettazione e creazione delle pagine del sito si sono utilizzate le linee guida del Corso di \textit{Tecnologie Web} e quelle del \textit{W3C}. Il codice è stato validato utilizzando il tool di validazione di 
	\href{https://validator.w3.org/}{W3C}.
	
Alcune delle regole più importanti sono:
\begin{itemize}
	\item Separazione della struttura, presentazione e contenuto: Non ci devono essere file script oppure fogli di stile nel codice \textit{HTML}. Queste devono essere scritte in file separati e poi si devono importare nell'\textit{header} delle pagine \textit{HTML}.
	\item \textit{Tag}: Tutti i \textit{tag} devono essere chiusi e si devono utilizzare i tag che migliorano l'accessibilità del sito.
	\item \textit{Metatag}: L’inserimento dei \textit{metatag} nella parte dello \textit{header} è un regola importantissima. Se inseriti correttamente migliorano l’accessibilità del sito verso i browser.Allo stesso tempo, l’utilizzo corretto delle \textit{keyword} aiuta in un \textit{ranking} migliore nella ricerca.
\end{itemize}

\subsection{PHP}
Il comportamento del sito lato server si gestisce da file PHP, i quali stabiliscono il comportamento generale delle pagine, interagiscono con il database e creano le sessioni di utilizzo per gli utenti loggati. I file creati sono:
\begin{itemize}
	\item db.php
	\item response\_manager.php
	\item test.php
\end{itemize}

\textbf{File dp.php}: In questo file ci sono tutte le classi e funzioni che servono per lo svolgimento delle \textit{query} del database e il controllo di eventuali errori di connessione o di interazione con esso. Abbiamo inizialmente la classe Constant dove vengono dichiarate le costanti che riguardano il database: \textit{(HOST\_DB, DATABASE\_NAME, USERNAME, PASSWORD)}.

Abbiamo poi la classe \textit{Service} la quale estende la classe \textit{Constant}. Questa classe serve per effettuare la connessione al database utilizzando le funzioni \textit{openConnection()} che restituisce un \textit{booleano} in base al risultato ottenuto e la funzione \textit{closeConnection()} per chiudere la connessione.

Queste funzioni verranno utilizzate poi in ogni pagina \textit{“namepagina.php”} dove serve stabilire una connessione in modo da poter effettuare poi una \textit{query} riguardante il risultato cercato.

L’altra parte del contenuto di \textit{db.php} sono le varie funzioni utilizzate per effettuare le \textit{query}. Alcune funzioni effettuano le \textit{query} che servono per le funzionalità che vengono utilizzate da un utente generico (non registrato), altre per le funzionalità disponibili per un utente registrato e infine quelle per le funzionalità dell'amministratore del sito.

\textbf{Utente generico:} Un utente generico ha la possibilità di visualizzare quasi tutte le pagine e le sezioni del sito, tranne alcune riservate per un utente registrato. Per fare questo si utilizzano diverse funzioni, tra le quale possiamo elencare:

\begin{itemize}
	\item \textbf{get\_book\_by\_(isbn/title)}: Queste funzioni creano la \textit{query} che effettua una ricerca secondo il parametro specificato.Restituiscono un messaggio di errore nel caso in cui qualcosa non è andato bene oppure quando non si è trovato nessun risultato. Utilizzati quando si effettua una ricerca nella barra di ricerca
	\item \textbf{get\_books\_by\_(author/genre)}: Se queste funzioni non trovano nessun risultato corrispondente al valore dato in \textit{input}, amplificano la ricerca cercando tutti i risultati che contengono quel valore nel campo cercato in qualsiasi posizione(utilizzando l’operatore \textbf{*campo cercato*} di \textit{SQL}). Utilizzabili quando si vogliono trovare i libri scritti da un certo autore oppure quando si vuole visualizzare la lista dei libri appartenenti a quel specifico genere.
	\item \textbf{get\_new\_books\_by\_genre}: Restituisce un \textit{array} contenente 5 nuovi libri aggiunti nella categoria cercata, secondo l’ordine decrescente riguardo la data di pubblicazione.
	\item \textbf{get\_genere\_by\_id}; Funzione che restituisce il genere con quel \textit{id}. Utilizzata quando si vuole visualizzare la pagina di un genere specifico dalla lista dei generi.
	\item \textbf{get\_bestsellers}: Funzione che esegue la \textit{query} la quale restituisce i libri \textit{bestseller} di tutti i generi in quel momento ordinati dal libro più venduto a quello meno venduto. Utilizzata nella pagina home e anche per la pagina Bestseller stessa.
	\item \textbf{signin}: Funzione che gestisce la creazione di un nuovo \textit{account} inserendolo nel \textit{database}. Nel caso si commettano errori(es. il tipo di uno dei valori in \textit{input} non è corretto) viene visualizzato un messaggio di errore.
	\item \textbf{get\_avg\_review}: Funzione che esegue la \textit{query} per il calcolo della media delle recensioni di un certo libro. Utilizzata nella pagina \textit{Libro} che contiene i dettagli di quel specifico libro.
	\item \textbf{get\_reviews\_by\_isbn}: Funzione che serve per trovare tutte le recensioni lasciate ad un libro. Si utilizza nella pagina \textit{Libro} nella parte dove vengono listati tutti le review.
	\item \textbf{get\_new\_books}: Funzione che esegue la \textit{query} che trova i libri che sono stati pubblicati recentemente. Necessaria per la parte \textit{“Nuove uscite”} della \textit{homepage}.
	\item \textbf{get\_books\_under\_5}: Funzione che esegue la \textit{query} che trova i libri che hanno prezzo minore di 5 euro. Necessaria per la parte “\textit{A meno di 5 euro”} nella \textit{homepage}.
	\item \textbf{get\_all\_books}: Funzione che restituisce la lista di tutti i libri presenti nel \textit{database}. Utilizzato nella pagina \textit{Libri} e inizialmente nella pagina \textit{Ricerca} (prima che venga effettuata una ricerca, vengono visualizzati alcuni libri del \textit{database} a destra della pagina, mentre a sinistra ci sono i filtri e la barra di ricerca).
	\item \textbf{get\_genres\_from\_isbn}: Questa funzione esegue la \textit{query} che trova i generi dove appartiene un libro con un certo \textit{isbn}. Viene utilizzata nella parte del \textit{"Dettagli Libro”} della pagina \textit{Libro} e anche nella parte dei filtri sui generi della pagina \textit{Ricerca}.
	\item \textbf{get\_books\_with\_offers}: Funzione che effettua la \textit{query} che selezioni i libri che hanno una riduzione di prezzo, ossia sono in offerta. Necessaria per la pagina \textit{Offerte}.
	\item \textbf{get\_active\_offer\_by\_isbn}: Funzione che trova quanto è lo sconto del prezzo di un specifico libro presenti nella pagina \textit{Offerte}. Necessario nella pagina \textit{Libro} dove viene visualizzato lo sconto attuale.
\end{itemize}

\textbf{Utente registrato:} Le funzioni necessarie per eseguire le funzionalità che un utente registrato possiede sono tutte quelle dell'utente generico(tranne la registrazione), con altre in più.
\begin{itemize}
	\item \textbf{login}: Funzione che utilizza la \textit{query} per ottenere i dati dei campi della tabella utente che appartengono all’indirizzo \textit{email} e \textit{password} inserite durante la fase di accesso.
	\item \textbf{get\_addresses}: Effettua una \textit{query} che trova tutti gli indirizzi appartenenti ad un id utente. Si utilizza nella sezione \textit{Indirizzi} di ogni \textit{account} dove si possono vedere gli indirizzi inseriti.
	\item \textbf{insert\_addresses}: Funzione che gestisce l’inserzione di un nuovo indirizzo nel \textit{database} collegato con l’\textit{account} del utente. Si utilizza nella parte \textit{“Inserisci indirizzo”} della pagina \textit{Indirizzi}.
	\item \textbf{insert\_review}: Funzione che gestisce l’inserzione di una nuova recensione da parte di un utente per uno libro specifico. Vengono visualizzati appositi messaggi di errore in caso qualcosa non sia andato bene.
	\item \textbf{get\_reviews\_by\_user/by\_user\_book}: La prima funzione estrae tutte le recensioni lasciate da quel utente per qualsiasi libro, mentre la seconda trova tutte le recensioni effettuate per un libro specifico.
	\item \textbf{insert\_order}: Funzione che aggiunge un nuovo ordine nella tabella \textit{Ordine} del \textit{database}. Vengono fatti vari controlli riguardo i campi dei dati inseriti e se si incontra alcun problema un apposito messaggio di errore viene visualizzato.
	\item \textbf{get\_orders}: ?
	\item \textbf{edit\_review/delete\_review}: Sono delle funzioni che utilizzano le \textit{query} per trovare le recensioni effettuate dall’utente e modificarle/cancellarle. Si utilizzano nella pagina \textit{Recensioni} che ogni utente ha dentro al proprio profilo. Se durante il processo si commettono errori, oppure qualcosa non va bene vengono visualizzati dei messaggi di errore.
	\item \textbf{restore\_code}: Funzione che inserisce nella tabella recupero la nuova coppia di valori del codice di conferma(\textit{id}) ricevuto e il codice identificativo dell’utente a cui appartiene. Nel caso in cui ci sia già un valore per questa \textit{tupla}, esso viene aggiornato.
	\item \textbf{restore\_psw}: Funzione che effettua una \textit{query} che cancella i vecchi valori dalla tabella recupero e aggiorna il nuovo \textit{password} dell'utente nella tabella utente. Viene utilizzata nei casi in cui l’utente ha dimenticato la \textit{password} e la recupera attraverso un codice ricevuto via \textit{mail}.
	\item \textbf{is\_code\_correct}: Funzione che esegue una \textit{query} che restituisce la \textit{tupla} nella tabella recupero con i corrispondenti valori inseriti per i campi id ed utente. Nel caso in cui non si trova un risultato, oppure succede un errore viene visualizzato un messaggio.
\end{itemize}

\textbf{Amministratore:} L’amministratore dispone delle funzionalità in più rispetto ai semplici utenti. Le funzioni che aiutano ad poter eseguire queste funzionalità sono:
\begin{itemize}
	\item \textbf{get\_utente\_by\_email/id}: Funzioni che effettuano la \textit{query} che ricerca l’utente a cui appartiene l’indirizzo \textit{email}/\textit{id} selezionato. Vengono utilizzate nella gestione degli utenti dalla parte del amministratore.
	\item \textbf{insert\_book}: Questa funzione permette di inserire un nuovo libro nel \textit{database}. Utilizzata nella pagina \textit{Aggiungi Libri}.
	\item \textbf{edit\_book}: Attraverso questa funzione è possibile modificare tutti oppure alcuni campi di un libro già presente nel \textit{database}. Utilizzata nella pagina \textit{Modifica Libro}.
\end{itemize}

\textbf{File response\_manager.php:} In questo file viene creata la classe \textit{response\_manager}. Sono dichiarate al suo interno variabili e metodi utilizzati per la gestione delle risposte e il controllo degli errori. Alcuni più importanti sono:
\begin{itemize}
	\item \textbf{\_\_construct}: che dichiara lo costruttore della classe \textit{response\_manager}
	\item \textbf{get\_result}: che prende il risultato della \textit{query} dal \textit{server MySQL }e lo trasporterà al \textit{server PHP}
	\item \textbf{set\_errorr\_messagge}: viene utilizzato per il controllo degli errori
\end{itemize}

\textbf{File NomePagina.php:} Sono i file che costruiscono le pagine del sito. Utilizzano le rispettive pagine \textit{HTML} che contengono l’effettivo \textit{HTML} e il \textit{file db.php}. Tra quelli che costruiscono e gestiscono le funzionalità più importanti del sito sono:

\end{document}